% SOURCE_AUTHORITY: sga5_fr_workpass.tex, lines 11935--11941 (LF-normalized unit).
% SOURCE_SHA256: 791F4EFFC5E02832D5D77ED03518C8156D6F07E4C8238B03545DB93D883FBB28
Um objeto $M$ de $\cC$ será chamado pseudocoerente se o for como complexo concentrado em grau zero; demonstra-se que, para isso, é necessário e suficiente que admita uma resolução à esquerda cujas componentes sejam objetos de $\cC_o$.

\begin{statement}{Proposição 5.1}
Um complexo $X$ é $\cC_o$-pseudocoerente se, e somente se, for isomorfo a um complexo limitado à direita cujas componentes são $\cC_o$-pseudocoerentes.
\end{statement}

A demonstração desta proposição figurará na primeira exposição do seminário sobre Riemann--Roch (SGA 6 I 2.7).
